\section{Optional, Non-Promoting Analyses}

The main evidence path is observational FCI followed by randomized assigned-arm
ITT. JCI is an exploratory analysis of randomized contexts that, when run,
preserves separate raw and constrained PAGs plus explicit assumption-linked
orientation deltas~\cite{mooij2020jci}. RFCI is an optional py-tetrad
sensitivity backend; its availability or failure does not block the no-Java
causal-learn path. Neither JCI nor RFCI may change the candidate universe,
selector ranks, bridge, assignment, outcomes, ITT, or evidence status.

Separately provenanced implementation markers, treatment fidelity, and
per-protocol summaries are optional diagnostics. They cannot promote an
evidence level, substitute for the independent Oracle, enter the primary PAG,
or support a mediation claim. The RQ4 expert study is separately governed by
its ethics decision, preregistration, case and randomization manifests, and
analysis plan; it cannot validate either observational prioritization or
randomized security effects.
